\begin{abstract}
Large language model (LLM) agents are increasingly augmented with tools,
persistent memory, retrieval components, and the ability to communicate with
other agents, which extends their capabilities but also broadens their attack
surface. This is particularly true for multi-agent systems, where sensitive
information can propagate through internal channels -- messages between
agents, tool arguments, memory writes, and logs -- before, or even without,
reaching the final user-facing output. This project proposes an
AutoResearch-driven red-teaming methodology for agentic and multi-agent LLM
systems, in which an LLM agent iteratively proposes attack variants, evaluates
them against a target system on a fixed task, scores the outcome, and retains
improving variants. We instantiate two independent optimization loops, one
maximizing synthetic sensitive-information exposure and the other maximizing
task-utility degradation and operational-cost amplification, and evaluate them
across AgentDojo and three multi-agent frameworks: AutoGen, CrewAI, and
LangGraph. By measuring both final-output and internal-channel leakage,
utility degradation, cost amplification, and cross-framework transferability,
we aim to determine whether automatically optimized attacks outperform random
and manually designed baselines, and whether attacks discovered on one
framework generalize to others.
\end{abstract}
